\QHeading{第1题　口算。}
\par\noindent


\begin{oralgrid}
\OralItem{1}{2.5\times 1000}{2500}
\OralItem{2}{30\div 100}{0.3}
\OralItem{3}{388\times 10}{3880}
\OralItem{4}{73.6\div 10}{7.36}
\OralItem{5}{0.8\times 100}{80}
\OralItem{6}{9\div 1000}{0.009}
\OralItem{7}{0.45\times 10}{4.5}
\OralItem{8}{67.8\div 100}{0.678}\end{oralgrid}
\par

\QHeading{口算思路补充（对应上表题号）}

\Answer{上面列出得数，下面按小题给出不同计算思路。计数单位法适用于小数，整数题可比较凑整、拆分或商不变等方法。}
\MethodTitle{第2小题：$30\div 100$}

\MethodBlock{口算思路一：同时去掉相同个数的0}{}{\begin{aligned}[t]&\phantom{={}}30\div 100\\
&=(30\div 10)\div (100\div 10)\\
&=3\div 10\\
&=0.3\end{aligned}}
\MethodTitle{第3小题：$388\times 10$}

\MethodBlock{口算思路一：凑整后补偿}{把接近整十、整百的因数凑整，再补回或减去多算的部分。}{\begin{aligned}[t]&\phantom{={}}388\times 10\\
&=10\times (390-2)\\
&=10\times 390-10\times 2\\
&=3900-20\\
&=3880\end{aligned}}

\MethodBlock{口算思路二：按数位拆分计算}{把一个因数拆成整十、整百数与余下部分之和。}{\begin{aligned}[t]&\phantom{={}}388\times 10\\
&=10\times (300+88)\\
&=10\times 300+10\times 88\\
&=3000+880\\
&=3880\end{aligned}}


\par\noindent

\QHeading{第2题　填一填。}

\par\noindent\makebox[\linewidth][c]{
\begin{tikzpicture}[x=1bp,y=-1bp,color=black,draw=black,text=black]
\begin{scope}\clip (56.0,229.0) rectangle (471.0,361.33);
\hdhead{50}{小数的意义与性质——相关链接}
\hdsection{1}{154}{口算。}
\q{99}{201}{180}{13}{$2.5\times 1000=$}
\q{99}{245}{180}{13}{$0.8\times 100=$}
\q{287}{201}{180}{13}{$30\div 100=$}
\q{287}{245}{180}{13}{$9\div 1000=$}
\q{478}{201}{180}{13}{$388\times 10=$}
\q{478}{245}{180}{13}{$0.45\times 10=$}
\q{674}{201}{180}{13}{$73.6\div 10=$}
\q{674}{245}{180}{13}{$67.8\div 100=$}
\hdsection{2}{300}{填一填。}
\q{108}{429}{60}{13}{48.3}
\q{254}{353}{85}{13}{4.83}
\begin{scope}[scale=.6666667,draw=black,text=black]\draw[draw=black,line width=.6pt] (156.000,439.000)--(242.000,370.000);\end{scope}
\begin{scope}[scale=.6666667,draw=black,text=black]\draw[draw=black,line width=.6pt] (156.000,439.000)--(248.000,413.000);\end{scope}
\begin{scope}[scale=.6666667,draw=black,text=black]\draw[draw=black,line width=.6pt] (156.000,439.000)--(250.000,452.000);\end{scope}
\begin{scope}[scale=.6666667,draw=black,text=black]\draw[draw=black,line width=.6pt] (156.000,439.000)--(253.000,492.000);\end{scope}
\begin{scope}[scale=.6666667,draw=black,text=black]\node[inner sep=0pt,font=\normalsize,text=black] at (201.000,383.000) {$\div 10$};\end{scope}
\begin{scope}[scale=.6666667,draw=black,text=black]\node[inner sep=0pt,font=\normalsize,text=black] at (215.000,409.000) {$\times 10$};\end{scope}
\begin{scope}[scale=.6666667,draw=black,text=black]\node[inner sep=0pt,font=\normalsize,text=black] at (226.000,438.000) {$\div 100$};\end{scope}
\begin{scope}[scale=.6666667,draw=black,text=black]\node[inner sep=0pt,font=\normalsize,text=black] at (228.000,470.000) {$\times 100$};\end{scope}
\q{244}{396}{28}{12}{（}
\q{313}{396}{25}{12}{）}
\q{244}{442}{28}{12}{（}
\q{313}{442}{25}{12}{）}
\q{244}{480}{28}{12}{（}
\q{313}{480}{25}{12}{）}
\q{495}{352}{94}{13}{0.729}
\q{510}{397}{94}{13}{72.9}
\q{503}{445}{94}{13}{7290}
\q{484}{484}{94}{13}{0.0729}
\q{665}{422}{94}{13}{7.29}
\begin{scope}[scale=.6666667,draw=black,text=black]\draw[draw=black,line width=.6pt] (569.000,367.000)--(658.000,432.000);\end{scope}
\begin{scope}[scale=.6666667,draw=black,text=black]\draw[draw=black,line width=.6pt] (562.000,408.000)--(658.000,432.000);\end{scope}
\begin{scope}[scale=.6666667,draw=black,text=black]\draw[draw=black,line width=.6pt] (560.000,452.000)--(658.000,432.000);\end{scope}
\begin{scope}[scale=.6666667,draw=black,text=black]\draw[draw=black,line width=.6pt] (572.000,490.000)--(658.000,432.000);\end{scope}
\begin{scope}[scale=.6666667,draw=black,text=black]\node[inner sep=0pt,font=\normalsize,text=black] at (632.000,376.000) {$\times 10$};\end{scope}
\begin{scope}[scale=.6666667,draw=black,text=black]\path[fill=white,draw=none] (561,387) rectangle (638,411);\end{scope}
\q{554}{388}{28}{11.5}{（}
\q{630}{388}{25}{11.5}{）}
\begin{scope}[scale=.6666667,draw=black,text=black]\path[fill=white,draw=none] (561,423) rectangle (638,447);\end{scope}
\q{554}{424}{28}{11.5}{（}
\q{630}{424}{25}{11.5}{）}
\begin{scope}[scale=.6666667,draw=black,text=black]\path[fill=white,draw=none] (561,457) rectangle (638,481);\end{scope}
\q{554}{458}{28}{11.5}{（}
\q{630}{458}{25}{11.5}{）}
\hdsection{3}{545}{算一算，填一填。}
\begin{scope}[scale=.6666667,draw=black,text=black]\draw[line width=.5pt] (100,605) rectangle (166,636);\end{scope}
\begin{scope}[scale=.6666667,draw=black,text=black]\draw[line width=.5pt] (262,605) rectangle (328,636);\end{scope}
\begin{scope}[scale=.6666667,draw=black,text=black]\draw[line width=.5pt] (429,605) rectangle (495,636);\end{scope}
\begin{scope}[scale=.6666667,draw=black,text=black]\draw[line width=.5pt] (596,605) rectangle (662,636);\end{scope}
\begin{scope}[scale=.6666667,draw=black,text=black]\node[inner sep=0pt,font=\normalsize,text=black] at (133.000,621.000) {8};\end{scope}
\begin{scope}[scale=.6666667,draw=black,text=black]\draw[draw=black,line width=.6pt,-{Stealth[length=4pt,width=3pt]}] (168.000,620.000)--(258.000,620.000);\end{scope}
\begin{scope}[scale=.6666667,draw=black,text=black]\node[inner sep=0pt,font=\normalsize,text=black] at (214.000,597.000) {$\div 10$};\end{scope}
\begin{scope}[scale=.6666667,draw=black,text=black]\draw[draw=black,line width=.6pt,-{Stealth[length=4pt,width=3pt]}] (330.000,620.000)--(420.000,620.000);\end{scope}
\begin{scope}[scale=.6666667,draw=black,text=black]\node[inner sep=0pt,font=\normalsize,text=black] at (376.000,597.000) {$\times 100$};\end{scope}
\begin{scope}[scale=.6666667,draw=black,text=black]\draw[draw=black,line width=.6pt,-{Stealth[length=4pt,width=3pt]}] (497.000,620.000)--(587.000,620.000);\end{scope}
\begin{scope}[scale=.6666667,draw=black,text=black]\node[inner sep=0pt,font=\normalsize,text=black] at (543.000,597.000) {$\div 1000$};\end{scope}
\begin{scope}[scale=.6666667,draw=black,text=black]\draw[line width=.5pt] (100,671) rectangle (166,702);\end{scope}
\begin{scope}[scale=.6666667,draw=black,text=black]\draw[line width=.5pt] (262,671) rectangle (328,702);\end{scope}
\begin{scope}[scale=.6666667,draw=black,text=black]\draw[line width=.5pt] (429,671) rectangle (495,702);\end{scope}
\begin{scope}[scale=.6666667,draw=black,text=black]\draw[line width=.5pt] (596,671) rectangle (662,702);\end{scope}
\begin{scope}[scale=.6666667,draw=black,text=black]\node[inner sep=0pt,font=\normalsize,text=black] at (133.000,687.000) {95};\end{scope}
\begin{scope}[scale=.6666667,draw=black,text=black]\draw[draw=black,line width=.6pt,-{Stealth[length=4pt,width=3pt]}] (168.000,686.000)--(258.000,686.000);\end{scope}
\begin{scope}[scale=.6666667,draw=black,text=black]\node[inner sep=0pt,font=\normalsize,text=black] at (214.000,663.000) {$\div 100$};\end{scope}
\begin{scope}[scale=.6666667,draw=black,text=black]\draw[draw=black,line width=.6pt,-{Stealth[length=4pt,width=3pt]}] (330.000,686.000)--(420.000,686.000);\end{scope}
\begin{scope}[scale=.6666667,draw=black,text=black]\node[inner sep=0pt,font=\normalsize,text=black] at (376.000,663.000) {$\times 10$};\end{scope}
\begin{scope}[scale=.6666667,draw=black,text=black]\draw[draw=black,line width=.6pt,-{Stealth[length=4pt,width=3pt]}] (497.000,686.000)--(587.000,686.000);\end{scope}
\begin{scope}[scale=.6666667,draw=black,text=black]\node[inner sep=0pt,font=\normalsize,text=black] at (543.000,663.000) {$\times 1000$};\end{scope}
\begin{scope}[scale=.6666667,draw=black,text=black]\draw[line width=.5pt] (100,738) rectangle (166,769);\end{scope}
\begin{scope}[scale=.6666667,draw=black,text=black]\draw[line width=.5pt] (262,738) rectangle (328,769);\end{scope}
\begin{scope}[scale=.6666667,draw=black,text=black]\draw[line width=.5pt] (429,738) rectangle (495,769);\end{scope}
\begin{scope}[scale=.6666667,draw=black,text=black]\draw[line width=.5pt] (596,738) rectangle (662,769);\end{scope}
\begin{scope}[scale=.6666667,draw=black,text=black]\node[inner sep=0pt,font=\normalsize,text=black] at (133.000,754.000) {8.5};\end{scope}
\begin{scope}[scale=.6666667,draw=black,text=black]\draw[draw=black,line width=.6pt,-{Stealth[length=4pt,width=3pt]}] (168.000,753.000)--(258.000,753.000);\end{scope}
\begin{scope}[scale=.6666667,draw=black,text=black]\node[inner sep=0pt,font=\normalsize,text=black] at (214.000,730.000) {$\times 100$};\end{scope}
\begin{scope}[scale=.6666667,draw=black,text=black]\draw[draw=black,line width=.6pt,-{Stealth[length=4pt,width=3pt]}] (330.000,753.000)--(420.000,753.000);\end{scope}
\begin{scope}[scale=.6666667,draw=black,text=black]\node[inner sep=0pt,font=\normalsize,text=black] at (376.000,730.000) {$\div 10$};\end{scope}
\begin{scope}[scale=.6666667,draw=black,text=black]\draw[draw=black,line width=.6pt,-{Stealth[length=4pt,width=3pt]}] (497.000,753.000)--(587.000,753.000);\end{scope}
\begin{scope}[scale=.6666667,draw=black,text=black]\node[inner sep=0pt,font=\normalsize,text=black] at (543.000,730.000) {$\div 1000$};\end{scope}
\hdsection{4}{803}{解决问题。}
\q{101}{851}{346}{12}{① 一种糖果，每千克售价12.8}
\q{101}{893}{346}{12}{　元。购买10千克、100千克}
\q{101}{935}{346}{12}{　分别需要多少钱?}
\q{485}{851}{350}{12}{② 把一个小数的小数点先向右}
\q{485}{893}{350}{12}{　移动两位，再向左移动一位，}
\q{485}{935}{350}{12}{　得到的数是36.7。原来的小}
\q{485}{977}{350}{12}{　数是多少?}

\begin{scope}[color=answerred,draw=answerred,text=answerred]\ans{145.33}{134.00}{60.00}{9.5}{$\vphantom{0}2500$}
\ans{262.67}{134.00}{60.00}{9.5}{$\vphantom{0}0.3$}
\ans{388.67}{134.00}{60.00}{9.5}{$\vphantom{0}3880$}
\ans{526.67}{134.00}{60.00}{9.5}{$\vphantom{0}7.36$}
\ans{138.00}{163.33}{60.00}{9.5}{$\vphantom{0}80$}
\ans{262.67}{163.33}{60.00}{9.5}{$\vphantom{0}0.009$}
\ans{396.67}{163.33}{60.00}{9.5}{$\vphantom{0}4.5$}
\ans{533.33}{163.33}{60.00}{9.5}{$\vphantom{0}0.678$}
\ans{180.00}{262.67}{33.33}{10}{$\vphantom{0}483$}
\ans{178.00}{294.00}{33.33}{10}{$\vphantom{0}0.483$}
\ans{176.67}{319.33}{33.33}{10}{$\vphantom{0}4830$}
\ans{383.33}{258.67}{33.33}{10}{$\vphantom{0}\div10$}
\ans{379.33}{282.67}{33.33}{10}{$\vphantom{0}\div1000$}
\ans{382.67}{305.33}{33.33}{10}{$\vphantom{0}\times100$}
\ans{186.67}{406.67}{33.33}{10}{$\vphantom{0}0.8$}
\ans{298.67}{406.00}{33.33}{10}{$\vphantom{0}80$}
\ans{408.67}{405.33}{33.33}{10}{$\vphantom{0}0.08$}
\ans{185.33}{450.67}{33.33}{10}{$\vphantom{0}0.95$}
\ans{298.00}{450.00}{33.33}{10}{$\vphantom{0}9.5$}
\ans{406.00}{450.00}{33.33}{10}{$\vphantom{0}9500$}
\ans{184.67}{495.33}{33.33}{10}{$\vphantom{0}850$}
\ans{298.67}{495.33}{33.33}{10}{$\vphantom{0}85$}
\ans{406.67}{494.67}{33.33}{10}{$\vphantom{0}0.085$}
\ans{69.33}{655.33}{222.67}{10}{10千克：$12.8\times10=128$（元）。\\答：购买10千克需要128元。\\100千克：$12.8\times100=1280$（元）。\\答：购买100千克需要1280元。}
\ans{326.67}{678.00}{220.00}{10}{按小数点移动顺序逆向还原：}
\ans{326.67}{693.33}{221.33}{10}{$\setlength{\jot}{-2pt}\begin{aligned}[t]&\phantom{={}}36.7\times10\div100\\&=367\div100\\&=3.67\end{aligned}$}
\ans{326.67}{741.33}{220.00}{10}{答：原来的小数是3.67。}
\end{scope}
\end{scope}\pgfresetboundingbox\path[use as bounding box] (56.0,229.0) rectangle (471.0,361.33);
\end{tikzpicture}}\par



\par



\par\noindent

\QHeading{第3题　算一算，填一填。}

\par\noindent\makebox[\linewidth][c]{
\begin{tikzpicture}[x=1bp,y=-1bp,color=black,draw=black,text=black]
\begin{scope}\clip (56.0,387.0) rectangle (447.0,533.33);
\hdhead{50}{小数的意义与性质——相关链接}
\hdsection{1}{154}{口算。}
\q{99}{201}{180}{13}{$2.5\times 1000=$}
\q{99}{245}{180}{13}{$0.8\times 100=$}
\q{287}{201}{180}{13}{$30\div 100=$}
\q{287}{245}{180}{13}{$9\div 1000=$}
\q{478}{201}{180}{13}{$388\times 10=$}
\q{478}{245}{180}{13}{$0.45\times 10=$}
\q{674}{201}{180}{13}{$73.6\div 10=$}
\q{674}{245}{180}{13}{$67.8\div 100=$}
\hdsection{2}{300}{填一填。}
\q{108}{429}{60}{13}{48.3}
\q{254}{353}{85}{13}{4.83}
\begin{scope}[scale=.6666667,draw=black,text=black]\draw[draw=black,line width=.6pt] (156.000,439.000)--(242.000,370.000);\end{scope}
\begin{scope}[scale=.6666667,draw=black,text=black]\draw[draw=black,line width=.6pt] (156.000,439.000)--(248.000,413.000);\end{scope}
\begin{scope}[scale=.6666667,draw=black,text=black]\draw[draw=black,line width=.6pt] (156.000,439.000)--(250.000,452.000);\end{scope}
\begin{scope}[scale=.6666667,draw=black,text=black]\draw[draw=black,line width=.6pt] (156.000,439.000)--(253.000,492.000);\end{scope}
\begin{scope}[scale=.6666667,draw=black,text=black]\node[inner sep=0pt,font=\normalsize,text=black] at (201.000,383.000) {$\div 10$};\end{scope}
\begin{scope}[scale=.6666667,draw=black,text=black]\node[inner sep=0pt,font=\normalsize,text=black] at (215.000,409.000) {$\times 10$};\end{scope}
\begin{scope}[scale=.6666667,draw=black,text=black]\node[inner sep=0pt,font=\normalsize,text=black] at (226.000,438.000) {$\div 100$};\end{scope}
\begin{scope}[scale=.6666667,draw=black,text=black]\node[inner sep=0pt,font=\normalsize,text=black] at (228.000,470.000) {$\times 100$};\end{scope}
\q{244}{396}{28}{12}{（}
\q{313}{396}{25}{12}{）}
\q{244}{442}{28}{12}{（}
\q{313}{442}{25}{12}{）}
\q{244}{480}{28}{12}{（}
\q{313}{480}{25}{12}{）}
\q{495}{352}{94}{13}{0.729}
\q{510}{397}{94}{13}{72.9}
\q{503}{445}{94}{13}{7290}
\q{484}{484}{94}{13}{0.0729}
\q{665}{422}{94}{13}{7.29}
\begin{scope}[scale=.6666667,draw=black,text=black]\draw[draw=black,line width=.6pt] (569.000,367.000)--(658.000,432.000);\end{scope}
\begin{scope}[scale=.6666667,draw=black,text=black]\draw[draw=black,line width=.6pt] (562.000,408.000)--(658.000,432.000);\end{scope}
\begin{scope}[scale=.6666667,draw=black,text=black]\draw[draw=black,line width=.6pt] (560.000,452.000)--(658.000,432.000);\end{scope}
\begin{scope}[scale=.6666667,draw=black,text=black]\draw[draw=black,line width=.6pt] (572.000,490.000)--(658.000,432.000);\end{scope}
\begin{scope}[scale=.6666667,draw=black,text=black]\node[inner sep=0pt,font=\normalsize,text=black] at (632.000,376.000) {$\times 10$};\end{scope}
\begin{scope}[scale=.6666667,draw=black,text=black]\path[fill=white,draw=none] (561,387) rectangle (638,411);\end{scope}
\q{554}{388}{28}{11.5}{（}
\q{630}{388}{25}{11.5}{）}
\begin{scope}[scale=.6666667,draw=black,text=black]\path[fill=white,draw=none] (561,423) rectangle (638,447);\end{scope}
\q{554}{424}{28}{11.5}{（}
\q{630}{424}{25}{11.5}{）}
\begin{scope}[scale=.6666667,draw=black,text=black]\path[fill=white,draw=none] (561,457) rectangle (638,481);\end{scope}
\q{554}{458}{28}{11.5}{（}
\q{630}{458}{25}{11.5}{）}
\hdsection{3}{545}{算一算，填一填。}
\begin{scope}[scale=.6666667,draw=black,text=black]\draw[line width=.5pt] (100,605) rectangle (166,636);\end{scope}
\begin{scope}[scale=.6666667,draw=black,text=black]\draw[line width=.5pt] (262,605) rectangle (328,636);\end{scope}
\begin{scope}[scale=.6666667,draw=black,text=black]\draw[line width=.5pt] (429,605) rectangle (495,636);\end{scope}
\begin{scope}[scale=.6666667,draw=black,text=black]\draw[line width=.5pt] (596,605) rectangle (662,636);\end{scope}
\begin{scope}[scale=.6666667,draw=black,text=black]\node[inner sep=0pt,font=\normalsize,text=black] at (133.000,621.000) {8};\end{scope}
\begin{scope}[scale=.6666667,draw=black,text=black]\draw[draw=black,line width=.6pt,-{Stealth[length=4pt,width=3pt]}] (168.000,620.000)--(258.000,620.000);\end{scope}
\begin{scope}[scale=.6666667,draw=black,text=black]\node[inner sep=0pt,font=\normalsize,text=black] at (214.000,597.000) {$\div 10$};\end{scope}
\begin{scope}[scale=.6666667,draw=black,text=black]\draw[draw=black,line width=.6pt,-{Stealth[length=4pt,width=3pt]}] (330.000,620.000)--(420.000,620.000);\end{scope}
\begin{scope}[scale=.6666667,draw=black,text=black]\node[inner sep=0pt,font=\normalsize,text=black] at (376.000,597.000) {$\times 100$};\end{scope}
\begin{scope}[scale=.6666667,draw=black,text=black]\draw[draw=black,line width=.6pt,-{Stealth[length=4pt,width=3pt]}] (497.000,620.000)--(587.000,620.000);\end{scope}
\begin{scope}[scale=.6666667,draw=black,text=black]\node[inner sep=0pt,font=\normalsize,text=black] at (543.000,597.000) {$\div 1000$};\end{scope}
\begin{scope}[scale=.6666667,draw=black,text=black]\draw[line width=.5pt] (100,671) rectangle (166,702);\end{scope}
\begin{scope}[scale=.6666667,draw=black,text=black]\draw[line width=.5pt] (262,671) rectangle (328,702);\end{scope}
\begin{scope}[scale=.6666667,draw=black,text=black]\draw[line width=.5pt] (429,671) rectangle (495,702);\end{scope}
\begin{scope}[scale=.6666667,draw=black,text=black]\draw[line width=.5pt] (596,671) rectangle (662,702);\end{scope}
\begin{scope}[scale=.6666667,draw=black,text=black]\node[inner sep=0pt,font=\normalsize,text=black] at (133.000,687.000) {95};\end{scope}
\begin{scope}[scale=.6666667,draw=black,text=black]\draw[draw=black,line width=.6pt,-{Stealth[length=4pt,width=3pt]}] (168.000,686.000)--(258.000,686.000);\end{scope}
\begin{scope}[scale=.6666667,draw=black,text=black]\node[inner sep=0pt,font=\normalsize,text=black] at (214.000,663.000) {$\div 100$};\end{scope}
\begin{scope}[scale=.6666667,draw=black,text=black]\draw[draw=black,line width=.6pt,-{Stealth[length=4pt,width=3pt]}] (330.000,686.000)--(420.000,686.000);\end{scope}
\begin{scope}[scale=.6666667,draw=black,text=black]\node[inner sep=0pt,font=\normalsize,text=black] at (376.000,663.000) {$\times 10$};\end{scope}
\begin{scope}[scale=.6666667,draw=black,text=black]\draw[draw=black,line width=.6pt,-{Stealth[length=4pt,width=3pt]}] (497.000,686.000)--(587.000,686.000);\end{scope}
\begin{scope}[scale=.6666667,draw=black,text=black]\node[inner sep=0pt,font=\normalsize,text=black] at (543.000,663.000) {$\times 1000$};\end{scope}
\begin{scope}[scale=.6666667,draw=black,text=black]\draw[line width=.5pt] (100,738) rectangle (166,769);\end{scope}
\begin{scope}[scale=.6666667,draw=black,text=black]\draw[line width=.5pt] (262,738) rectangle (328,769);\end{scope}
\begin{scope}[scale=.6666667,draw=black,text=black]\draw[line width=.5pt] (429,738) rectangle (495,769);\end{scope}
\begin{scope}[scale=.6666667,draw=black,text=black]\draw[line width=.5pt] (596,738) rectangle (662,769);\end{scope}
\begin{scope}[scale=.6666667,draw=black,text=black]\node[inner sep=0pt,font=\normalsize,text=black] at (133.000,754.000) {8.5};\end{scope}
\begin{scope}[scale=.6666667,draw=black,text=black]\draw[draw=black,line width=.6pt,-{Stealth[length=4pt,width=3pt]}] (168.000,753.000)--(258.000,753.000);\end{scope}
\begin{scope}[scale=.6666667,draw=black,text=black]\node[inner sep=0pt,font=\normalsize,text=black] at (214.000,730.000) {$\times 100$};\end{scope}
\begin{scope}[scale=.6666667,draw=black,text=black]\draw[draw=black,line width=.6pt,-{Stealth[length=4pt,width=3pt]}] (330.000,753.000)--(420.000,753.000);\end{scope}
\begin{scope}[scale=.6666667,draw=black,text=black]\node[inner sep=0pt,font=\normalsize,text=black] at (376.000,730.000) {$\div 10$};\end{scope}
\begin{scope}[scale=.6666667,draw=black,text=black]\draw[draw=black,line width=.6pt,-{Stealth[length=4pt,width=3pt]}] (497.000,753.000)--(587.000,753.000);\end{scope}
\begin{scope}[scale=.6666667,draw=black,text=black]\node[inner sep=0pt,font=\normalsize,text=black] at (543.000,730.000) {$\div 1000$};\end{scope}
\hdsection{4}{803}{解决问题。}
\q{101}{851}{346}{12}{① 一种糖果，每千克售价12.8}
\q{101}{893}{346}{12}{　元。购买10千克、100千克}
\q{101}{935}{346}{12}{　分别需要多少钱?}
\q{485}{851}{350}{12}{② 把一个小数的小数点先向右}
\q{485}{893}{350}{12}{　移动两位，再向左移动一位，}
\q{485}{935}{350}{12}{　得到的数是36.7。原来的小}
\q{485}{977}{350}{12}{　数是多少?}

\begin{scope}[color=answerred,draw=answerred,text=answerred]\ans{145.33}{134.00}{60.00}{9.5}{$\vphantom{0}2500$}
\ans{262.67}{134.00}{60.00}{9.5}{$\vphantom{0}0.3$}
\ans{388.67}{134.00}{60.00}{9.5}{$\vphantom{0}3880$}
\ans{526.67}{134.00}{60.00}{9.5}{$\vphantom{0}7.36$}
\ans{138.00}{163.33}{60.00}{9.5}{$\vphantom{0}80$}
\ans{262.67}{163.33}{60.00}{9.5}{$\vphantom{0}0.009$}
\ans{396.67}{163.33}{60.00}{9.5}{$\vphantom{0}4.5$}
\ans{533.33}{163.33}{60.00}{9.5}{$\vphantom{0}0.678$}
\ans{180.00}{262.67}{33.33}{10}{$\vphantom{0}483$}
\ans{178.00}{294.00}{33.33}{10}{$\vphantom{0}0.483$}
\ans{176.67}{319.33}{33.33}{10}{$\vphantom{0}4830$}
\ans{383.33}{258.67}{33.33}{10}{$\vphantom{0}\div10$}
\ans{379.33}{282.67}{33.33}{10}{$\vphantom{0}\div1000$}
\ans{382.67}{305.33}{33.33}{10}{$\vphantom{0}\times100$}
\ans{186.67}{406.67}{33.33}{10}{$\vphantom{0}0.8$}
\ans{298.67}{406.00}{33.33}{10}{$\vphantom{0}80$}
\ans{408.67}{405.33}{33.33}{10}{$\vphantom{0}0.08$}
\ans{185.33}{450.67}{33.33}{10}{$\vphantom{0}0.95$}
\ans{298.00}{450.00}{33.33}{10}{$\vphantom{0}9.5$}
\ans{406.00}{450.00}{33.33}{10}{$\vphantom{0}9500$}
\ans{184.67}{495.33}{33.33}{10}{$\vphantom{0}850$}
\ans{298.67}{495.33}{33.33}{10}{$\vphantom{0}85$}
\ans{406.67}{494.67}{33.33}{10}{$\vphantom{0}0.085$}
\ans{69.33}{655.33}{222.67}{10}{10千克：$12.8\times10=128$（元）。\\答：购买10千克需要128元。\\100千克：$12.8\times100=1280$（元）。\\答：购买100千克需要1280元。}
\ans{326.67}{678.00}{220.00}{10}{按小数点移动顺序逆向还原：}
\ans{326.67}{693.33}{221.33}{10}{$\setlength{\jot}{-2pt}\begin{aligned}[t]&\phantom{={}}36.7\times10\div100\\&=367\div100\\&=3.67\end{aligned}$}
\ans{326.67}{741.33}{220.00}{10}{答：原来的小数是3.67。}
\end{scope}
\end{scope}\pgfresetboundingbox\path[use as bounding box] (56.0,387.0) rectangle (447.0,533.33);
\end{tikzpicture}}\par



\par

\par\noindent

\QHeading{第4题　解决问题。}

\Problem{① 一种糖果，每千克售价12.8元。购买10千克、100千克分别需要多少钱?}
\par


\MethodTitle{方法一}
\Answer{10千克：$12.8\times10=128$（元）。\par 答：购买10千克需要128元。\par 100千克：$12.8\times100=1280$（元）。\par 答：购买100千克需要1280元。}

\MethodTitle{方法二：先求10千克的钱，再扩大到100千克}
\Answer{100千克是10千克的10倍。}
\Answer{\textbf{10千克糖果的价钱}\par \Formula{\begin{aligned}[t]&\phantom{={}}12.8\times 10\\
&=128\text{（元）}\end{aligned}}}
\Answer{\textbf{100千克糖果的价钱}\par \Formula{\begin{aligned}[t]&\phantom{={}}128\times 10\\
&=1280\text{（元）}\end{aligned}}}
\Answer{答：购买10千克需要128元。\par 答：购买100千克需要1280元。}
\Problem{② 把一个小数的小数点先向右移动两位，再向左移动一位，得到的数是36.7。原来的小数是多少?}

\MethodTitle{方法一}
\Answer{按小数点移动顺序逆向还原：}
\Answer{\Formula{\begin{aligned}[t]&\phantom{={}}36.7\times10\div100\\&=367\div100\\&=3.67\end{aligned}}}
\Answer{答：原来的小数是3.67。}

\MethodTitle{方法二：先合并小数点移动效果}
\Answer{先右移两位再左移一位，总效果是向右移一位，得到的数是原数的10倍。}
\Answer{\Formula{\begin{aligned}[t]&\phantom{={}}36.7\div 10\\
&=3.67\end{aligned}}}
\Answer{答：原来的小数是3.67。}
